\subsection*{Experimental setup and frequency transfer}

The comparison links two optical lattice clocks at different institutions:
an ytterbium (Yb) clock at IAPMST in Wuhan and a strontium (Sr) clock at
USTC in Shanghai, about $\siteKm$\,km apart. The two sites are connected by
a phase-stabilised fibre link of $\fibreKm$\,km single-pass length
($\roundKm$\,km loop-back), operated at \qty{1550}{nm}. The route is
divided into spans of tens of kilometres, each terminated by a relay
station that carries a narrow-linewidth laser together with an
erbium-doped fibre amplifier and a fibre-noise-cancellation loop; the
added instability of the link is kept below the clock noise level for
averaging times beyond $100$\,s.

Both ends use independent optical frequency combs to bridge the
spectral gap between the clock transitions and the telecom carrier. On the
Shanghai side the Sr \qty{698}{nm} transition is transferred via a
\qty{1397}{nm} continuous laser to the \qty{1550}{nm} carrier; on the
Wuhan side the received carrier is linked to the Yb \qty{578}{nm}
transition through the \qty{1156}{nm} light. All local frequency
references and phase-locked loops are derived from the comb repetition
rate $f_{\mathrm{rep}} = \fRepMHz$\,MHz. The observable is the
out-of-loop beat between the received carrier and the local
Yb-referenced comb, converted to fractional units with the transfer
frequency $F_{1550} = \fFifteenTHz$\,THz (see Appendix~\ref{app:chain}).

\subsection*{Observation campaign and per-segment data}

The campaign spanned \nDays~days, from 29 June to 26 August 2026, in
\nStages~measurement stages. Synchronous-operation windows were first
defined from the lock status of the two clocks and of the beat signals;
within each window the longest continuous valid record was retained.
The selection used lock status only and did not refer to the resulting
ratio deviations. In total \nSegAll~jump-free segments were recorded,
carrying \nSamplesAll~valid seconds.

The ninth recorded segment is excluded from the primary analysis. Its Sr
systematic frequency-shift parameter sits between the values of the two
operating periods and its provenance has not been confirmed by the
experiment team; the segment is therefore held pending verification. The
exclusion is based on parameter provenance, not on the resulting ratio
deviation, and its effect is quantified in the sensitivity check
(Section~\ref{sec:results-sensitivity}). The primary data set thus
comprises \nSeg~segments and \nSamples~valid seconds
(\nHours~effective hours), i.e.\ \nSegNine~seconds fewer than the full
record.

\subsection*{Clock-ratio determination and statistical model}

Each retained segment yields a clock-ratio estimate $R_i$ from the mean
of its retained beat samples, inverted through the fixed frequency-chain
relation (Appendix~\ref{app:chain}) and corrected for the static
gravitational redshift $\delta_g = \deltaGV$ derived from the levelled
geopotential difference (see below). The
per-segment statistical uncertainty $u_i$ is obtained from the
overlapping Allan deviation of the beat fractional frequency: a
$\sigma_y \propto \tau^{-1/2}$ fit over the range
$128~\mathrm{s} \le \tau \le 0.25\,T_i$, extrapolated to the segment
duration $T_i$.

The per-segment uncertainties do not fully explain the observed scatter
between segments, so the ratio is combined with four estimators:
fixed-effect weighted least squares (WLS, weights $1/u_i^2$), Birge
scaling, the Mandel--Paule random-effect model, and a Bayesian
random-effect model in which the excess scatter $\xi$ is a free
parameter whose posterior uncertainty is propagated into the final
result. The Bayesian random-effect centre is adopted as the primary
statistic; the others are reported as cross-checks. Five of the
\nSeg~segments (groups \modelLimitedGroups) have stability fits that
fall outside the white-frequency-noise band and are carried with a
model-limited flag; this is one motivation for the random-effect
treatment. The full model definitions are given in
Appendix~\ref{app:stats}.

\subsection*{Levelling determination of the geopotential difference}

The static gravitational-redshift correction is fixed by the
geopotential difference between the two clocks' atomic transition
positions, determined by first-order spirit levelling together with
simultaneous gravimetry. A levelling line of about $\levelKm$\,km was
run between the Wuhan benchmark WS-03 (at IAPMST) and the Shanghai
benchmark WS-01 (at USTC) through Ezhou, Huangshi, Huanggang, Lu'an,
Hefei, Wuhu, Guangde and Suzhou, using digital levels with sight
lengths below \qty{30}{m} and gravity observed at every levelling point
with two relative gravimeters. The levelled normal heights were
converted to orthometric heights for the geopotential integration,
$\dW_{AB} = -\sum_i g_i\,\Delta H_i$, where $g_i$ is the mean gravity
and $\Delta H_i$ the orthometric height difference of the $i$-th
segment. The resulting geopotential difference between the two atomic
transitions is $\dW = \dWVal \pm \dWUnc~\mathrm{m^2\,s^{-2}}$. The
per-kilometre misclosure of the levelling is $\levelMm$\,mm, below the
first-order limit.

\subsection*{Uncertainty budget}

Table~\ref{tab:budget} lists the uncertainty components of the
Yb/Sr ratio. The statistical term is the posterior standard deviation
of the Bayesian random-effect model applied to the \nSeg-segment data.
The Sr and Yb systematic terms are the published evaluations of the two
clocks~\citep{jia2026sr, zhu2026yb}; the value adopted for the Yb clock
resolves an internal-versus-published discrepancy in favour of the
published evaluation. The static-potential term follows from the
levelling uncertainty. The fibre-link and frequency-comb terms are
experiment-side upper bounds. The time-varying
gravitational-redshift term is carried as an uncertainty item,
quantified from the scale over which the segment-averaged shift varies
across the campaign, and is not applied as a correction to the reported
central value. The components are treated as independent and combined
in quadrature.
